%% LyX 2.5.1 created this file.  For more info, see https://www.lyx.org/.
%% Do not edit unless you really know what you are doing.
\documentclass[10pt,english,10pt]{beamer}
\usepackage{lmodern}
\usepackage[T1]{fontenc}
\usepackage[utf8]{inputenc}
\setcounter{tocdepth}{1}
\usepackage{babel}
\usepackage[authoryear]{natbib}
\PassOptionsToPackage{normalem}{ulem}
\usepackage{ulem}
\ifx\hypersetup\undefined
  \AtBeginDocument{%
    \hypersetup{}
  }
\else
  \hypersetup{}
\fi

\makeatletter
%%%%%%%%%%%%%%%%%%%%%%%%%%%%%% Textclass specific LaTeX commands.
% this default might be overridden by plain title style
\newcommand\makebeamertitle{\frame{\maketitle}}%
% (ERT) argument for the TOC
\AtBeginDocument{%
  \let\origtableofcontents=\tableofcontents
  \def\tableofcontents{\@ifnextchar[{\origtableofcontents}{\gobbletableofcontents}}
  \def\gobbletableofcontents#1{\origtableofcontents}
}

%%%%%%%%%%%%%%%%%%%%%%%%%%%%%% User specified LaTeX commands.



\usepackage{tikz}
\usetikzlibrary{positioning}
\usepackage{appendixnumberbeamer}

\usepackage{graphicx}
\usepackage{subfig}

\usetheme[progressbar=frametitle,block=fill,subsectionpage=progressbar]{metropolis}

% margin
\setbeamersize{text margin right=1.5cm}

% colors
\colorlet{DarkRed}{red!70!black}
\setbeamercolor{normal text}{fg=black}
\setbeamercolor{alerted text}{fg=DarkRed}
\setbeamercolor{progress bar}{fg=DarkRed}
\setbeamercolor{button}{bg=DarkRed}

% width of seperators
\makeatletter
\setlength{\metropolis@titleseparator@linewidth}{1pt}
\setlength{\metropolis@progressonsectionpage@linewidth}{1pt}
\setlength{\metropolis@progressinheadfoot@linewidth}{1pt}
\makeatother

% new alert block
\newlength\origleftmargini
\setlength\origleftmargini\leftmargini
\setbeamertemplate{itemize/enumerate body begin}{\setlength{\leftmargini}{4mm}}
\let\oldalertblock\alertblock
\let\oldendalertblock\endalertblock
\def\alertblock{\begingroup \setbeamertemplate{itemize/enumerate body begin}{\setlength{\leftmargini}{\origleftmargini}} \oldalertblock}
\def\endalertblock{\oldendalertblock \endgroup}
\setbeamertemplate{mini frame}{}
\setbeamertemplate{mini frame in current section}{}
\setbeamertemplate{mini frame in current subsection}{}
\setbeamercolor{section in head/foot}{fg=normal text.bg, bg=structure.fg}
\setbeamercolor{subsection in head/foot}{fg=normal text.bg, bg=structure.fg}

% footer
\makeatletter
\setbeamertemplate{footline}{%
    \begin{beamercolorbox}[colsep=1.5pt]{upper separation line head}
    \end{beamercolorbox}
    \begin{beamercolorbox}{section in head/foot}
      \vskip1pt\insertsectionnavigationhorizontal{\paperwidth}{}{\hskip0pt plus1filll \insertframenumber{} / \inserttotalframenumber \hskip2pt}\vskip3pt%
    \end{beamercolorbox}%
    \begin{beamercolorbox}[colsep=1.5pt]{lower separation line head}
    \end{beamercolorbox}
}
\makeatother

% toc
\setbeamertemplate{section in toc}{\hspace*{1em}\inserttocsectionnumber.~\inserttocsection\par}
\setbeamertemplate{subsection in toc}{\hspace*{2em}\inserttocsectionnumber.\inserttocsubsectionnumber.~\inserttocsubsection\par}

% code
\usepackage{xcolor}
\usepackage{listings}

\definecolor{codegray}{rgb}{0.5,0.5,0.5}
\definecolor{background}{HTML}{F5F5F5}
\definecolor{keyword}{HTML}{4B69C6}
\definecolor{string}{HTML}{448C27}
\definecolor{comment}{HTML}{448C27}

\usepackage{inconsolata}
\lstdefinestyle{mystyle}{
    commentstyle=\color{comment},
    keywordstyle=\color{keyword},
    stringstyle=\color{string},
    basicstyle=\ttfamily,
    breakatwhitespace=false,
    breaklines=true,
    captionpos=b,
    keepspaces=true,
    numbersep=5pt,
    showspaces=false,
    showstringspaces=false,
    showtabs=false,
    tabsize=4,
	showlines=true
}

\lstset{style=mystyle}
% Added by lyx2lyx
\setlength{\parskip}{\smallskipamount}
\setlength{\parindent}{0pt}

\makeatother

\begin{document}
\title{1. Introduction\vspace{-2mm}}
\subtitle{Adv. Macro: Heterogenous Agent Models}
\author{Raphaël Huleux}
\date{2026}

{
\setbeamertemplate{footline}{}
\begin{frame}
\maketitle
\end{frame}
}

\addtocounter{framenumber}{-1}

\begin{frame}<beamer>
\frametitle{Plan}
\tableofcontents[]
\end{frame}

\section{Introduction}
\begin{frame}{Introduction}
\vspace{-2mm}
\begin{itemize}
\item <+->\textbf{Teacher: }Raphaël Huleux
\item <+->\textbf{Central economic questions:}\vspace{1mm}
\begin{enumerate}
\item <+->What role does heterogeneity play for understanding consumption-saving
dynamics in partial equilibrium?\\
(incl. single-agent problems)\vspace{1mm}
\item <+->What role does heterogeneity play for understanding\\
secular trends in general equilibrium?\vspace{1mm}
\item <+->What explains the level and dynamics of heterogeneity/inequality?\vspace{1mm}
\item <+->What role does heterogeneity play for understanding business-cycle
fluctuations in general equilibrium?\vspace{1mm}
\end{enumerate}
\item <+->\textbf{Central technical method:} Programming in Python

\textbf{Prerequisite:} \emph{Intro. to Programming and Numerical Analysis}

\textbf{Complicated:} \emph{Close to the research frontier}
\end{itemize}
\end{frame}
%
\begin{frame}{Macroeconomic Models with Heterogeneous Agents}
\begin{itemize}
\item <+->\vspace{-2mm}\textbf{Model components:}
\begin{enumerate}
\item Optimizing individual agents (households + firms)
\item Idiosyncratic and aggregate risk
\item Information flows (who knows what when $\Rightarrow$ often everything)
\item Market clearing (Walras vs. search-and-match)
\end{enumerate}
\item <+->\textbf{Insurance/markets:}

\emph{Complete} $\rightarrow$ idiosyncratic risk insured away $\sim$
representative agent

\emph{Incomplete} $\rightarrow$ agents need to \emph{self-insure}
\item <+->\textbf{Heterogeneity:}

\emph{Ex ante }in\emph{ }preferences, abilities etc.

\emph{Ex post} after realization of idiosyncratic shocks
\item <+->\textbf{HANC: }Heterogeneous Agent \emph{Neo-Classical }model\\
{\footnotesize (Aiyagari-Bewley-Hugget-Imrohoroglu or Standard Incomplete
Markets model)}{\footnotesize\par}
\item <+->\textbf{HANK: }Heterogeneous Agent \emph{New Keynesian} model\\
(i.e. include price and wage setting frictions)
\end{itemize}
\end{frame}
%
\begin{frame}{History of heterogeneous agent macro}
\begin{enumerate}
\item Heathcote et al. (2009), >>Quantitative Macroeconomics with Heterogeneous
Households<<
\item Kaplan and Violante (2018), >>Microeconomic Heterogeneity and Macroeconomic
Shocks<<
\item Cherrier et al. (2023), >>Household Heterogeneity in Macroeconomic
Models: A Historical Perspective<<
\item Auclert et. al. (2025), >>Fiscal and Monetary Policy with Heterogeneous
Agents<<
\end{enumerate}
\end{frame}
%

\section{Structure}
\begin{frame}{Teaching method}
\begin{itemize}
\item <+->\textbf{Lectures: }Tuesday 12-15

\textasciitilde 2 hours of >>normal<< lecture

\textasciitilde 1 hour of code review and problem solving (no exercise
classes)
\item <+->\textbf{Content:}
\begin{enumerate}
\item Explanation of computational methods
\item Exposition of central stylized results
\item Code examples for central methods and mechanisms
\item Discussion of research papers
\end{enumerate}
\item <+->\textbf{Material:}\vspace{1mm}

{\small Web: }\textcolor{DarkRed}{{\small\href{https://sites.google.com/view/numeconcph-advmacrohet/}{sites.google.com/view/numeconcph-advmacrohet/}}}{\small\par}

{\small Git: }\textcolor{DarkRed}{{\small\href{https://github.com/NumEconCopenhagen/AdvMacroHet}{github.com/numeconcopenhagen/adv-macro-het}}}{\small\par}
\item <+->\textbf{Code:}
\begin{enumerate}
\item We provide code you will build upon
\item Based on the \textcolor{DarkRed}{\href{https://github.com/NumEconCopenhagen/GEModelTools}{GEModelTools}}
package
\end{enumerate}
\end{itemize}
\end{frame}
%
\begin{frame}{Assignments and exam}
\begin{itemize}
\item <+->Individual \textbf{assignments }(hand-in on Absalon)
\begin{enumerate}
\item <+->\textbf{Assignment I}

\uline{Deadline}: 23rd of October (\emph{must be approved before exam})
\item <+->\textbf{Assignment II}

\uline{Deadline}: 20th of November (\emph{must be approved before
exam})
\item <+->\textbf{Assignment III} with essay on a relevant model extension
of own choice and a simple implementation

\uline{Deadline for proposal}: 10th of December

\uline{Deadline for peer feedback}: 14th of December (\emph{exam requirement})
\end{enumerate}
\item <+->All\textbf{ feedback }can be used to improve assignments before
the exam
\item <+->\textbf{Exam:}
\begin{enumerate}
\item Hand-in\textbf{ $3\times$assignments}
\item \textbf{36 hour take-home:} Programming of new extension \\
+ analysis of model + interpretation of results
\end{enumerate}
\end{itemize}
\end{frame}
%
\begin{frame}{Python}
\begin{enumerate}
\item \textbf{Assumed knowledge:} From \textcolor{DarkRed}{\href{https://sites.google.com/view/numeconcph-introprog/home}{Introduction to Programming and Numerical Analysis}}
you are assumed to know the basics of
\begin{enumerate}
\item Python
\item VSCode
\item git
\end{enumerate}
\item \textbf{Updated Python:} Install (or re-install) newest Anaconda
\item \textbf{Packages:}\texttt{ }{\small\texttt{pip install quantecon,
EconModel, consav}}{\small\par}
\item \textbf{GEMoodel tools:}
\begin{enumerate}
\item Clone the \texttt{GEModelTools} repository
\item Locate repository in command prompt
\item Run \texttt{pip install -e .}
\end{enumerate}
\end{enumerate}
\end{frame}
%
\begin{frame}{Course plan}
\begin{itemize}
\item \textbf{Lecture 1-2:} Consumption-saving models
\item \textbf{Lecture 3-6:} General equilibrium model\\
(stationary equilibrium + transition path)
\item \textbf{Lecture 7:} Wealth inequality
\item \textbf{Lecture 8-13:} HANK
\item \textbf{Lecture 14:} Exam and perspectives
\end{itemize}
\end{frame}
%

\section{Learning goals}
\begin{frame}{Knowledge}
\begin{enumerate}
\item {\small Account for, formulate and interpret precautionary saving models }{\small\par}
\item {\small Account for stochastic and non-stochastic simulation methods }{\small\par}
\item {\small Account for, formulate and interpret general equilibrium models
with ex ante and ex post heterogeneity, idiosyncratic and aggregate
risk, and with and without pricing frictions }{\small\par}
\item {\small Discuss the difference between the stationary equilibrium,
the transition path and the dynamic equilibrium }{\small\par}
\item {\small Discuss the relationship between various equilibrium concepts
and their solution methods}{\small\par}
\item {\small Identify and account for methods for analyzing the dynamic
distributional effects of long-run policy (e.g. taxation and social
security) and short-run policy (e.g. monetary and fiscal policy)}{\small\par}
\end{enumerate}
\end{frame}
%
\begin{frame}{Skills}
\begin{enumerate}
\item {\small Solve precautionary saving problems with dynamic programming
and simulate behavior with stochastic and non-stochastic techniques }{\small\par}
\item {\small Solve general equilibrium models with ex ante and ex post heterogeneity,
idiosyncratic and aggregate risk, and with and without pricing frictions
(stationary equilibrium, transition path, dynamic equilibrium) }{\small\par}
\item {\small Analyze dynamics of income and wealth inequality }{\small\par}
\item {\small Analyze transitional and permanent structural changes (e.g.
inequality trends and the long-run decline in the interest rate) }{\small\par}
\item {\small Analyze the dynamic distributional effects of long-run policy
(e.g. taxation and social security) and short-run policy (e.g. monetary
and fiscal policy)}{\small\par}
\end{enumerate}
\end{frame}
%
\begin{frame}{Competencies}
\begin{enumerate}
\item {\small Independently formulate, discuss and assess research on both
the causes and effects of heterogeneity and risk for both long-run
and short-run outcomes }{\small\par}
\item {\small Discuss and assess the importance of how heterogeneity and
risk is modeled for questions about both long-run and short-run dynamics}{\small\par}
\end{enumerate}
\end{frame}
%

\section{Programming in Python}
\begin{frame}[fragile]{Classes}
\begin{itemize}
\item A \textbf{class} defines the \textbf{type} of an object
\begin{itemize}
\item \lstinline{.attribute}, state
\item \lstinline{.method()}, action (incl. changing self)
\end{itemize}
\item \textbf{Inheritance} (of methods) (\lstinline{class Child(Parent)})
\end{itemize}
\begin{lstlisting}[language=Python,basicstyle=\linespread{1.3}\ttfamily\footnotesize,
numbers=left,frame=single,backgroundcolor=\color{background}]
class Parent:
    def __init__(self,value): self.value = value
    def double(self): return self.value * 2

class Child(Parent):
    def half(self): return self.value / 2

child = Child(10)
print(child.value) # 10
print(child.double()) # 20
print(child.half()) # 5.0
\end{lstlisting}
\end{frame}
%
\begin{frame}[fragile]{References}
\begin{itemize}
\item \textbf{Variables are references to an instance of an object}
\item \lstinline{=} \textbf{assigns a reference} (\emph{not a copy!})
\end{itemize}
\textbf{Question:} What does \lstinline{a} end up as? What if \lstinline{a = [1,2,3]}?

\begin{lstlisting}[language=Python,basicstyle=\linespread{1.3}\ttfamily\footnotesize,
numbers=left,frame=single,backgroundcolor=\color{background}]
a = np.array([1,2,3])
b = a
c = a[1:] # slicing
b[0] = 3 # indexing
c[0] = 3
\end{lstlisting}
\end{frame}
%
\begin{frame}[fragile]{Types and in-place operations}
\begin{itemize}
\item \textbf{Atomic types:} \lstinline{int, float, str, bool,} etc.
\item \textbf{Containers }\lstinline{list, tuple, dict, set, np.array,}
etc.
\item \textbf{Mutables} (e.g. \lstinline{list, np.array})\textbf{ }can
change in-place
\begin{enumerate}
\item \textbf{In-place operators} (\lstinline{+=, -=} etc.)
\item \textbf{Slicing}: \lstinline{x[:] = x + y}
\end{enumerate}
\item \textbf{Immutables} (e.g. atomic types and tuples) can never change
\end{itemize}
\textbf{Questions: }What does \lstinline{y} end up as?

\begin{lstlisting}[language=Python,basicstyle=\linespread{1.3}\ttfamily\footnotesize,
numbers=left,frame=single,backgroundcolor=\color{background}]
x = np.array([1,2,3])
y = x
x += 1
x[:] = x + 1
x = x + 1
\end{lstlisting}
\end{frame}
%
\begin{frame}[fragile]{Functions and scope}
\begin{itemize}
\item \textbf{Functions are objects} (can e.g. be arguments in functions)

Unlike in math:
\begin{enumerate}
\item Can change its arguments (side-effects)
\item Can call itself (recursion)
\end{enumerate}
\item Decorators change function behavior (e.g. @numba.njit)
\item Variables can both be \textbf{local scope }(good) or \textbf{global
scope} (bad)
\end{itemize}
\textbf{Questions: }What is the output?

\begin{lstlisting}[language=Python,basicstyle=\linespread{1.3}\ttfamily\footnotesize,
numbers=left,frame=single,backgroundcolor=\color{background}]
a = 1
def f(x):
	return x+a
print(f(1))
a = 2
print(f(1))
\end{lstlisting}
\end{frame}
%
\begin{frame}[fragile]{Conditionals and loops}
\begin{itemize}
\item \textbf{Comparison }(\lstinline{==, !=, <, <=, not, and, or} etc.)
\item \textbf{Conditionals }(\lstinline{if, elif, else})
\item \textbf{Loops} (\lstinline{for, while, continue, break})
\item \textbf{Convergence} (tolerance in optimizer or root-finder/equation-solver)
\end{itemize}
\textbf{Questions: }How could this be implemented with a \lstinline{while}
loop?

\begin{lstlisting}[language=Python,basicstyle=\linespread{1.3}\ttfamily\footnotesize,
numbers=left,frame=single,backgroundcolor=\color{background}]
x = x0
for i in range(n):
    y = evaluate(x)
    if check(y): break
    x = update(x,y)
else:
    raise ValueError('did not converge')
\end{lstlisting}
\end{frame}
%
\begin{frame}[fragile]{Decimal numbers are not exact}
\begin{itemize}
\item \textbf{Never use exactness for decimal numbers }
\begin{itemize}
\item Order of computation matter
\item Best with numbers are around 1 (underflow and overflow)
\end{itemize}
\item Division, exp, log etc. are (costly) approximations
\item \textbf{Function approximation and interpolation often needed}
\end{itemize}
\textbf{Questions: }Which are \lstinline{True} and which are \lstinline{False}?

\begin{lstlisting}[language=Python,basicstyle=\linespread{1.3}\ttfamily\footnotesize,
numbers=left,frame=single,backgroundcolor=\color{background}]
print(0.1 + 0.2 == 0.3)
print(0.5 + 0.5 == 1.0)
print(np.isclose(0.1+0.2,0.3))
print(np.isclose(1e-200*1e200*1e200*1e-200,1.0))
print(np.isinf(1e-200*(1e200*1e200)*1e-200))
print(np.isclose(1e200*(1e-200*1e-200)*1e200,0.0))
\end{lstlisting}
\end{frame}
%
\begin{frame}[fragile]{Pseudo random numbers}
\begin{itemize}
\item \textbf{Only one seed }(randomness not assured across seeds)
\item State of random number generator can be reset
\item \textbf{Monte Carlo} simulation and integration
\begin{enumerate}
\item Static alternative: Use \textbf{quadrature rules }
\item Dynamic alternative: Discretize and derive \textbf{transition matrix}
\end{enumerate}
\end{itemize}
\textbf{Questions: }What is \lstinline{z} equal to?

\begin{lstlisting}[language=Python,basicstyle=\linespread{1.3}\ttfamily\footnotesize,
numbers=left,frame=single,backgroundcolor=\color{background}]
rng = np.random.default_rng(123)
s = rng.bit_generator.state
x = rng.normal(size=5)
y = rng.normal(size=5)
rng.bit_generator.state = s
z = rng.normal(size=5)
\end{lstlisting}
\end{frame}
%
\begin{frame}{Documentation and debugging}
\begin{itemize}
\item \textbf{No code is self-explanatory} (for others, incl. future you)
\item \textbf{Write documentation} (use \emph{github-copilot})
\begin{enumerate}
\item The comments explain humans what the code does.
\item The code makes the computer do what the comments say
\end{enumerate}
\item Important \textbf{design patterns}:
\begin{enumerate}
\item Use namespaces (be aware of scope) and meaningful names
\item No repetition of code-lines $\Rightarrow$ single-purpose functions/methods
\item Use \lstinline{assert} (also print and plot intermediate results)
\item Use \lstinline{try-except}
\end{enumerate}
\item \textbf{Run from top to bottom} (make \uline{\href{https://sites.google.com/view/numeconcph-introprog/guides/vscode}{shortcut}})

Replication: \uline{\href{https://datacodestandard.org/}{datacodestandard.org}}
\item \textbf{Debugging} (see 02. Debugging.ipynb)
\begin{enumerate}
\item Errors are (almost) always simple
\item Go through code step-by-step (\emph{manually} or \emph{debugger})
\end{enumerate}
\end{itemize}
\end{frame}
%
\begin{frame}{Numba and EconModelClass}
\begin{itemize}
\item \textbf{Numba:} Faster code (03. Numba.ipynb)
\item \textbf{EconModel} (04. EconModelClass.ipynb)
\begin{enumerate}
\item Make it easy to write well-structured code
\item Provide standard functionality for copying, saving and loading
\item Provide an easy interface to call numba JIT compilled functions
\item Provide an easy interface to call C++ functions\\
(not relevant in this course)
\end{enumerate}
\end{itemize}
\end{frame}
%

\section{Consumption-Saving}
\begin{frame}{Consumption-saving}

\emph{We start on the slides for lecture 2}
\end{frame}
%

\end{document}
